\begin{figure}[!htbp]
\centering
\begin{tikzpicture}[
  x=1cm,y=1cm,
  every node/.style={font=\small},
  retained/.style={draw,rounded corners=2pt,fill=black!4,
    minimum height=0.68cm,inner xsep=8pt},
  omitted/.style={draw=black!70,dashed,rounded corners=2pt,
    text=black!85,minimum height=0.68cm,inner xsep=7pt},
  link/.style={-{Stealth[length=1.7mm]},semithick},
  discarded/.style={link,dashed,draw=black!70}
]
\node (f) at (0,0) {$F_n$};
\node[retained] (sig) at (3.1,0)
  {$\sigma_1<\cdots<\sigma_i$};
\node[omitted] (tail) at (7.25,0) {discard $\le N_-+1$};
\node (onea) at (10.1,0) {$1$};
\draw[link] (f) -- (sig);
\draw[discarded] (sig) -- (tail);
\draw[discarded] (tail) -- (onea);

\node (oneb) at (0,-1.2) {$1$};
\node[omitted] (head) at (2.85,-1.2) {discard $\le N_++1$};
\node[retained] (tau) at (7,-1.2)
  {$\tau_j<\cdots<\tau_t$};
\node (b) at (10.1,-1.2) {$b_n$};
\draw[discarded] (oneb) -- (head);
\draw[discarded] (head) -- (tau);
\draw[link] (tau) -- (b);

\node[align=center] at (5.05,-2.48) {
  $Y\setminus U\subset\sigma_i^+\,,\qquad
   U\cap V=\varnothing\,,\qquad
   Y\setminus V\subset\tau_j^-$\\[5pt]
  $\sigma_i^+\cup\tau_j^-=Y\quad\Longrightarrow\quad
   \sigma_i<\tau_j$
};
\node[retained,inner xsep=12pt] (joined) at (5.05,-3.82)
  {$\sigma_1<\cdots<\sigma_i<\tau_j<\cdots<\tau_t$};
\node (fj) at (0,-3.82) {$F_n$};
\node (bj) at (10.1,-3.82) {$b_n$};
\draw[link] (fj) -- (joined);
\draw[link] (joined) -- (bj);
\end{tikzpicture}
\caption{Concatenating annular chains after bounded pruning. Choose
fixed disjoint neighborhoods $U$ of $q$ and $V$ of $z$, where
$F_n$ shrinks to $q$ and $b_n\to z\ne q$.
The terminal positive side of the first retained chain contains
$Y\setminus U$; the initial negative side of the second contains
$Y\setminus V$. Their union is $Y$, which is exactly the annulus-order
criterion for joining the chains. The diagram records order and
containment only; the sides need not be connected.}
\label{fig:annular-chains}
\end{figure}
